% Preliminary SRS + Traceability Matrix
% Project: Digital Procurement & PLS Seedbed (permissioned blockchain)
% LaTeX source intended for compilation with pdflatex / lualatex / Overleaf
\documentclass[11pt,a4paper]{article}
\usepackage[margin=1in]{geometry}
\usepackage{longtable}
\usepackage{booktabs}
\usepackage{hyperref}
\usepackage{array}
\usepackage{enumitem}
\usepackage{tabularx}
\usepackage{caption}
\usepackage{fancyhdr}
\pagestyle{fancy}
\fancyhead[L]{Digital Procurement \& PLS Seedbed}
\fancyhead[R]{Preliminary SRS}
\fancyfoot[C]{\thepage}
\title{Preliminary Software Requirements Specification (SRS) \\and Traceability Matrix\\\large Digital Procurement + Permissioned Blockchain + SME Market Access + PLS Seedbed}
\author{Prepared by: ChatGPT (preliminary public-source draft)}
\date{\today}
\begin{document}
\maketitle
\begin{abstract}
This document is a preliminary Software Requirements Specification (SRS) and traceability matrix for a digital procurement platform built on a permissioned blockchain, intended to support SME market access and seedbed financing using Shariah-compliant profit–loss sharing (PLS) instruments.\
\noindent
This draft is derived solely from public-domain sources (standards, vendor docs, regulatory guidance, whitepapers, academic literature, and open-source POCs). Each requirement is tagged with the source category and suggested stakeholder roles; stakeholder validation is required to finalize requirements.
\end{abstract}
\tableofcontents
\clearpage
\section{Document purpose and scope}
This SRS provides a structured set of candidate functional and non-functional requirements, compliance constraints, and a traceability matrix mapping requirements to public sources and assumed stakeholders. The system scope includes: \begin{itemize}[noitemsep]
\item SME onboarding and identity (DID/VC-enabled)
\item Procure-to-pay lifecycle (tender, PO, delivery, invoice, settlement)
\item Permissioned blockchain ledger and smart-contract templates
\item PLS product support (musharakah/mudharabah-like templates) and tokenisation of receivables
\item Integrations: banks (payments, financing), ERPs, logistics/3PL, standards (EPCIS/GS1, ISO20022)
\item Governance, audit, Shariah review, and regulatory reporting
\end{itemize}

\section{Assumed stakeholder roles}
(Used to tag requirements in traceability matrix)
\begin{itemize}
\item SMEs / Suppliers
\item Buyers / Procurement Officers (public \& private)
\item Participating Banks (PLS providers)
\item Regulators (financial, procurement, data protection) and Shariah Boards
\item Network / Consortium Admins and Operators
\item Logistics / 3PL Providers
\item Auditors / Compliance Officers
\item Developers / Integrators (API consumers)
\end{itemize}

\section{Definitions, acronyms, and conventions}
\begin{description}
\item[PLS] Profit-Loss Sharing (musharakah / mudharabah style)
\item[DID] Decentralized Identifier (W3C)
\item[VC] Verifiable Credential
\item[EPCIS] GS1 event capture standard for supply chain visibility
\item[SLA] Service Level Agreement
\item[KYC/AML] Know-Your-Customer / Anti-Money Laundering
\end{description}

\section{Project constraints and implementation assumptions}
\label{sec:constraints}

The initial implementation of the Digital Procurement \& PLS Seedbed is subject to practical delivery, budget, and architecture constraints. These constraints materially affect scope selection, technology choices, and MVP definition.

\subsection{Primary project constraints}
\begin{itemize}
\item \textbf{Schedule constraint.} The MVP must be designed, developed, tested, and demonstrated within \textbf{12 weeks (3 months)}. Therefore, only features that are implementable within this timebox shall be included in the first release.
\item \textbf{Cost constraint.} The solution must be \textbf{low-cost to deploy and operate} during the MVP stage. High upfront infrastructure, specialist platform administration, and expensive integration patterns shall be avoided unless strictly necessary for the pilot objectives.
\item \textbf{Resource constraint.} The project assumes a small development team with limited time for platform engineering, DevOps hardening, and enterprise-grade distributed network administration.
\end{itemize}

\subsection{Derived architectural constraints}
The following constraints are derived from the schedule and cost limits above:

\begin{itemize}
\item \textbf{No full-scale permissioned blockchain rollout in MVP.} A full consortium-grade Hyperledger Fabric deployment with multiple independently managed nodes, certificate authorities, ordering services, channel governance, and production-grade operations is considered too costly and too time-consuming for the initial 12-week MVP.
\item \textbf{Preference for centralized application core.} The MVP shall prioritize a \textbf{centralized platform architecture} using a conventional relational or document-oriented database, standard backend services, and workflow automation for core business processes.
\item \textbf{Hybrid architecture permitted if blockchain is retained.} If blockchain is still required for research or demonstrative purposes, the MVP shall adopt a \textbf{hybrid architecture}: core transaction processing, workflow, and application state remain off-chain in a centralized platform, while selected transaction hashes or verification proofs may be anchored to a permissioned blockchain.
\item \textbf{Blockchain use must be minimal and justified.} Any blockchain component included in the MVP must have a narrowly defined purpose such as integrity verification, audit proof anchoring, or research validation. Blockchain shall not be used for all business data or all application logic during the MVP unless a clear benefit is demonstrated.
\item \textbf{ERP integration shall be limited in MVP scope.} Direct integration with ERP or accounting systems is recognized as a major cost and complexity driver. Therefore, the MVP should use either mock APIs, CSV import/export, or a limited proof-of-concept integration rather than full production-grade ERP connectivity.
\item \textbf{Node hosting and operational overhead must be minimized.} The MVP shall avoid architecture choices that require extensive node hosting, multi-organization network maintenance, or complex blockchain governance operations during the first release.
\item \textbf{Preference for standard components.} The system should use commodity infrastructure and widely available components, including standard databases, API-based integration, and workflow automation mechanisms, to reduce setup cost and shorten delivery time.
\end{itemize}

\subsection{MVP boundary implications}
Given the above constraints, the MVP should focus on demonstrating the essential business flow rather than full technical decentralization. The first release should prioritize:
\begin{itemize}
\item supplier onboarding,
\item procure-to-pay workflow tracking,
\item basic auditability,
\item simple financing or PLS workflow simulation,
\item lightweight API integration patterns,
\item optional blockchain-backed verification only where necessary.
\end{itemize}

Features that require heavy infrastructure, multi-party production governance, or extensive enterprise integration may be deferred to later phases.

\subsection{Technology decision implication}
For the MVP, a \textbf{centralized-first} or \textbf{hybrid centralized + blockchain anchoring} design is the preferred baseline. A full Hyperledger Fabric implementation should be treated as a later-phase enhancement unless stakeholder validation explicitly confirms sufficient budget, time, and operational capacity.

\section{High-level system context and architecture}
The architecture described below must satisfy the project design constraints defined in Section~\ref{sec:design_constraints}.
\begin{itemize}
\item Baseline: permissioned, consortium blockchain (e.g., Hyperledger Fabric / Besu) with off-chain components for heavy data, analytics, and reporting.
\item Pattern: hybrid microservices + event-driven messaging + on-chain chaincode for core settlement/escrow/PLS logic. EPCIS-compatible event capture layer for logistics.
\end{itemize}

\section{Functional requirements}
\label{sec:functional}
Requirements are prefixed Rxx. Each entry contains: brief requirement, category, source category (public documents), stakeholders, and short acceptance criteria.
\begin{enumerate}[label=R\arabic*:,leftmargin=*,itemsep=6pt]
\item \textbf{R01 — DID-based decentralized identity for SMEs.} \
\textit{Type:} Functional / Data. \textit{Source category:} Standards (W3C DID/VC). \textit{Stakeholders:} SMEs, Network Admins. \textit{Acceptance:} SME can register and present a verifiable credential; verification returns valid cryptographic proof.
\item \textbf{R02 — KYC/AML onboarding workflow.} \
\textit{Type:} Functional / Compliance. \textit{Source:} Regulatory (AML/KYC). \textit{Stakeholders:} Banks, Regulators. \textit{Acceptance:} All transacting SMEs must pass KYC checks or be flagged; audit log retained per jurisdictional retention rules.
\item \textbf{R03 — Permissioned network membership \& role management.} \
\textit{Type:} Functional / Security. \textit{Source:} Platform docs (Hyperledger Fabric). \textit{Stakeholders:} Consortium Admins. \textit{Acceptance:} Admins can add/remove members; access control enforced by channels/ACLs.
\item \textbf{R04 — EPCIS-compatible event recording for shipment/receipt.} \
\textit{Type:} Functional / Interoperability. \textit{Source:} Standards (GS1 EPCIS). \textit{Stakeholders:} Buyers, Logistics. \textit{Acceptance:} Events ingest/export validate against EPCIS schema.
\item \textbf{R05 — Immutable ledger audit trail for procure-to-pay events.} \
\textit{Type:} Functional / Compliance. \textit{Source:} Industry pilots (World Bank, WEF). \textit{Stakeholders:} Auditors, Regulators. \textit{Acceptance:} Each PO→Delivery→Invoice→Settlement step recorded with timestamp \& signer hash.
\item \textbf{R06 — Smart-contract template for order acceptance \& escrow.} \
\textit{Type:} Functional. \textit{Source:} Vendor/POC docs. \textit{Stakeholders:} Buyers, SMEs, Banks. \textit{Acceptance:} Smart contract enforces PO terms and escrow flows; funds moved and released upon proof conditions.
\item \textbf{R07 — PLS contract model (musharakah/mudharabah).} \
\textit{Type:} Functional / Shariah compliance. \textit{Source:} Regulatory + AAOIFI/IFSB guidance. \textit{Stakeholders:} Banks, Shariah Board. \textit{Acceptance:} PLS contract deployable with parameters (profit ratio, loss allocation); ledger records distribution per contract.
\item \textbf{R08 — Tokenized representation of receivables for financing.} \
\textit{Type:} Functional / Integration. \textit{Source:} Tokenization whitepapers, industry reports. \textit{Stakeholders:} Banks, Financiers. \textit{Acceptance:} Invoice token can be issued, transferred, and redeemed; token provenance traces to original invoice.
\item \textbf{R09 — Endorsement policy \& validation for transactions.} \
\textit{Type:} Security / Platform. \textit{Source:} Hyperledger Fabric docs. \textit{Stakeholders:} Network Admins. \textit{Acceptance:} Transactions accepted only with required endorsers' signatures; others rejected.
\item \textbf{R10 — Privacy channels / private data collections.} \
\textit{Type:} Security / Architecture. \textit{Source:} Fabric private data / channels docs. \textit{Stakeholders:} Buyers, Banks. \textit{Acceptance:} Sensitive payment terms visible only to permitted parties; others see hashed proofs.
\item \textbf{R11 — API gateway with ISO20022-compatible payment messaging.} \
\textit{Type:} Integration / Non-functional. \textit{Source:} ISO20022 + vendor guides. \textit{Stakeholders:} Banks, Integrators. \textit{Acceptance:} Payment messages map to ISO20022 fields and are accepted by bank sandbox(s).
\item \textbf{R12 — Event-driven architecture for async reconciliation.} \
\textit{Type:} Non-functional / Design. \textit{Source:} Industry whitepapers (Deloitte/Oracle). \textit{Stakeholders:} Developers, Operators. \textit{Acceptance:} Reconciliation completes within agreed SLA after event processing; failed events retried.
\item \textbf{R13 — Performance baseline (configurable tx/sec).} \
\textit{Type:} Non-functional / Performance. \textit{Source:} Platform benchmarks. \textit{Stakeholders:} Network Ops. \textit{Acceptance:} System sustains configured tx/sec under load-testing with acceptable latency.
\item \textbf{R14 — Data governance \& machine-readable procurement metadata.} \
\textit{Type:} Data / Compliance. \textit{Source:} OECD digital procurement guidance. \textit{Stakeholders:} Procurement Officers, Regulators. \textit{Acceptance:} Tenders published with machine-readable metadata per recommended schema.
\item \textbf{R15 — Audit reporting \& export for regulators.} \
\textit{Type:} Functional / Compliance. \textit{Source:} OECD / World Bank pilots. \textit{Stakeholders:} Regulators. \textit{Acceptance:} Regulator can request \& receive signed ledger export bundle within SLA.
\item \textbf{R16 — Versioning \& upgradeable smart-contract patterns.} \
\textit{Type:} Maintainability / Functional. \textit{Source:} Platform best practices. \textit{Stakeholders:} Developers, Shariah Board. \textit{Acceptance:} Contract upgrades are auditable and backward-compatible or provide migration scripts.
\item \textbf{R17 — Role-based UI \& dashboards.} \
\textit{Type:} Functional / UX. \textit{Source:} Industry pilots. \textit{Stakeholders:} Buyers, Banks, Admins. \textit{Acceptance:} Each role sees tailored dashboard and can perform role-specific actions.
\item \textbf{R18 — On-chain evidence for delivery (IoT/QR proofs).} \
\textit{Type:} Integration / Functional. \textit{Source:} GS1 / IoT pilots. \textit{Stakeholders:} Logistics, Buyers. \textit{Acceptance:} Delivery event recorded and verifiable via signed IoT/QR proof.
\item \textbf{R19 — Data retention \& legal-hold policy.} \
\textit{Type:} Compliance. \textit{Source:} Regulatory guidance. \textit{Stakeholders:} Auditors, Legal. \textit{Acceptance:} Records retained per jurisdiction; legal-hold prevents deletion.
\item \textbf{R20 — Shariah governance workflow for PLS approvals.} \
\textit{Type:} Functional / Compliance. \textit{Source:} Malaysia SC, AAOIFI. \textit{Stakeholders:} Shariah Board, Banks. \textit{Acceptance:} Each PLS product passes Shariah checklist and is recorded before activation.
\item \textbf{R21 — Dispute resolution \& on-chain arbitration record.} \
\textit{Type:} Governance. \textit{Source:} WEF governance guidance. \textit{Stakeholders:} Buyers, SMEs, Auditors. \textit{Acceptance:} Dispute lifecycle recorded on-chain with evidence attachments and final status.
\item \textbf{R22 — Access logging \& cryptographic non-repudiation.} \
\textit{Type:} Security / Audit. \textit{Source:} NIST / ISO27001. \textit{Stakeholders:} Auditors, Network Ops. \textit{Acceptance:} All accesses logged with cryptographic signatures enabling non-repudiation.
\item \textbf{R23 — Interoperability adapter for legacy ERPs.} \
\textit{Type:} Integration. \textit{Source:} Vendor integration guides. \textit{Stakeholders:} Buyers, SMEs. \textit{Acceptance:} Adapter mapping passes sample PO/invoice integration tests.
\item \textbf{R24 — Consent management \& data privacy controls.} \
\textit{Type:} Privacy. \textit{Source:} OECD privacy guidance. \textit{Stakeholders:} SMEs, Regulators. \textit{Acceptance:} SMEs can grant/revoke consent; exports honor consent state.
\item \textbf{R25 — Monitoring, SLAs \& incident management.} \
\textit{Type:} Ops. \textit{Source:} Vendor ops guides. \textit{Stakeholders:} Network Ops, Admins. \textit{Acceptance:} Alerts for node failure and incident runbooks trigger notifications.
\item \textbf{R26 — Test harness with reproducible PLS scenarios.} \
\textit{Type:} Verification. \textit{Source:} Open-source POCs. \textit{Stakeholders:} Developers, QA, Banks. \textit{Acceptance:} Test suite reproduces profit/loss calculations across sample scenarios.
\item \textbf{R27 — Identity \& credential revocation mechanisms.} \
\textit{Type:} Security. \textit{Source:} W3C DID/VC revocation patterns. \textit{Stakeholders:} Network Admins. \textit{Acceptance:} Revoke API invalidates credential verification.
\item \textbf{R28 — Regulatory sandbox mode \& audit hooks.} \
\textit{Type:} Compliance / Functional. \textit{Source:} World Bank / national sandbox guidance. \textit{Stakeholders:} Regulators, Banks. \textit{Acceptance:} Sandbox mode enables enhanced logs and restricted deployment scope.
\item \textbf{R29 — Multi-jurisdiction configuration (retention/consent).} \
\textit{Type:} Compliance / Non-functional. \textit{Source:} OECD + national examples. \textit{Stakeholders:} Legal, Admins. \textit{Acceptance:} Configurable rules applied per jurisdiction in tests.
\item \textbf{R30 — Governance model: consortium bylaws \& on-chain policy registry.} \
\textit{Type:} Governance. \textit{Source:} WEF/World Bank governance docs. \textit{Stakeholders:} Consortium Admins, Regulators. \textit{Acceptance:} Policy changes require quorum and are recorded on-chain.
\end{enumerate}

\section{Non-functional requirements (selected)}
\begin{enumerate}[label=NF\arabic*:,leftmargin=*,itemsep=6pt]
\item \textbf{NF01 — Security baseline.} System follows ISO27001-aligned controls and NIST recommendations for identity, key management, and logging.
\item \textbf{NF02 — Performance / scalability.} Baseline configurable target e.g., 200 tx/sec for initial consortium; scalable to higher throughput with sharding/ordering improvements.
\item \textbf{NF03 — Availability.} 99.9\% availability target for core services during business hours; disaster recovery plan with RTO/RPO defined.
\item \textbf{NF04 — Privacy.} Payment terms and commercial-sensitive data stored in private collections or off-chain vaults; hashes and proofs stored on chain.
\item \textbf{NF05 — Maintainability / Upgrades.} Smart-contract upgrade path defined, migration scripts and test harness for any on-chain migration.
\end{enumerate}

\section{Compliance \& governance requirements}
List of high-risk regulatory constraints (examples):
\begin{itemize}
\item Shariah compliance for PLS products — profit sharing, loss allocation rules, audit trails.
\item AML/KYC and CDD obligations for financial participants.
\item Data protection \& retention laws per jurisdiction.
\item Public procurement transparency rules (machine-readable tenders per OECD guidance).
\end{itemize}

\section{Data model \& interfaces (high level)}
\begin{itemize}
\item Identity: DID + Verifiable Credential (W3C) for SMEs and orgs.
\item Event model: EPCIS event envelope for shipment/delivery/invoice events.
\item Payment messages: ISO20022-compatible fields for settlement \& reconciliation.
\item Smart-contract interface: chaincode functions for PO creation, accept/fulfil, escrow, issue-token, transfer-token, settle-PLS.
\end{itemize}

\section{Design constraints}
\label{sec:design_constraints}

The following constraints restrict the architecture, implementation approach, and scope of the MVP system. These constraints arise from project delivery limitations, infrastructure cost considerations, and stakeholder feedback regarding operational complexity.

\begin{enumerate}[label=C\arabic*:,leftmargin=*]
\item \textbf{C01 — Timebox constraint.} The MVP shall be deliverable within 12 calendar weeks, including development, testing, and demonstration.

\item \textbf{C02 — Low-cost deployment constraint.} The MVP shall use infrastructure and deployment patterns that minimize upfront and recurring operational cost.

\item \textbf{C03 — Minimal operational complexity constraint.} The MVP shall avoid solutions requiring substantial node hosting, distributed ledger administration, or complex consortium governance in the first release.

\item \textbf{C04 — Limited integration constraint.} ERP/accounting integration in the MVP shall be limited to lightweight API integration, mock services, or file-based exchange unless explicitly prioritized by stakeholders.

\item \textbf{C05 — Centralized-first architecture constraint.} The default MVP architecture shall be centralized unless blockchain functionality is proven necessary for a specific requirement.

\item \textbf{C06 — Hybrid blockchain constraint.} If blockchain is included in the MVP, it shall be limited to integrity verification, audit anchoring, or selective proof recording, while operational data and workflows remain primarily off-chain.

\item \textbf{C07 — Scope minimization constraint.} Requirements that significantly increase implementation time, hosting cost, or integration complexity may be deferred from the MVP to later system phases.
\end{enumerate}

\section{Traceability matrix}
The traceability matrix below maps the functional requirements (R01..R30) to public source categories and assumed stakeholders. This is a compact extract; the full matrix with direct source links is provided in Appendix A (bibliography).

\begin{longtable}{p{0.08\textwidth} p{0.35\textwidth} p{0.23\textwidth} p{0.23\textwidth}}
\caption{Traceability matrix (compact)}\\
\toprule
Req ID & Short requirement & Source category & Stakeholders \\
\midrule
\endfirsthead
\toprule
Req ID & Short requirement & Source category & Stakeholders \\
\midrule
\endhead
R01 & DID-based SME identity & W3C DID/VC (standards) & SMEs, Network Admins \\
R02 & KYC/AML onboarding workflow & Regulatory AML/KYC guidance & Banks, Regulators \\
R03 & Permissioned membership management & Hyperledger Fabric docs & Consortium Admins \\
R04 & EPCIS-compatible event recording & GS1 EPCIS & Buyers, Logistics \\
R05 & Immutable ledger audit trail & WEF / World Bank pilots & Auditors, Regulators \\
R06 & Smart-contract escrow & Vendor POCs & Buyers, SMEs, Banks \\
R07 & PLS contract model & Malaysia SC; AAOIFI/IFSB & Banks, Shariah Board \\
R08 & Tokenized receivables & Tokenisation whitepapers & Banks, Financiers \\
R09 & Endorsement policy & Fabric docs & Network Admins \\
R10 & Privacy collections & Fabric private data & Buyers, Banks \\
R11 & ISO20022 payment API & ISO20022 docs & Banks, Integrators \\
R12 & Event-driven reconciliation & Deloitte/Oracle whitepapers & Developers, Ops \\
R13 & Performance baseline & Platform benchmarks & Network Ops \\
R14 & Machine-readable procurement metadata & OECD guidance & Procurement Officers, Regulators \\
R15 & Audit export for regulators & OECD / World Bank & Regulators, Auditors \\
R16 & Versioning smart-contracts & Platform best practices & Developers, Shariah Board \\
R17 & Role-based UI & Industry pilots (WFP/IBM) & Buyers, Banks \\
R18 & On-chain IoT proofs & GS1 / IoT pilots & Logistics, Buyers \\
R19 & Retention & Regulatory guidance & Legal, Auditors \\
R20 & Shariah governance workflow & Malaysia SC, AAOIFI & Shariah Board, Banks \\
R21 & Dispute resolution record & WEF governance docs & Buyers, SMEs \\
R22 & Access logging & NIST / ISO27001 & Auditors, Ops \\
R23 & ERP adapter & Vendor integration guides & Buyers, SMEs \\
R24 & Consent management & OECD privacy guidance & SMEs, Regulators \\
R25 & Monitoring & Vendor ops docs & Network Ops \\
R26 & PLS test harness & Open-source POCs & Developers, Banks \\
R27 & Credential revocation & W3C DID/VC patterns & Network Admins \\
R28 & Regulatory sandbox mode & World Bank / national sandboxes & Regulators, Banks \\
R29 & Multi-jurisdiction support & OECD + national examples & Legal, Admins \\
R30 & Consortium governance & WEF / World Bank governance & Consortium Admins \\
\bottomrule
\end{longtable}

\section{Validation gaps and risk register (high level)}
\begin{itemize}
\item Business rules for PLS splits and exceptions require Shariah board confirmation (high risk).\
\item KYC thresholds and acceptable identity evidence for SMEs (medium risk).\
\item Jurisdictional retention/consent conflicts (multi-jurisdiction) (high risk).\
\item Throughput targets vs chosen platform capability (performance risk).\
\end{itemize}

\section{Appendix A: Selected bibliography (public sources used)}
Below is a curated set of public-domain sources used to derive the candidate requirements. For each, consult the document for technical details and canonical field lists.
\begin{enumerate}
\item OECD. "Digital Transformation of Public Procurement" (2025 report).
\item World Bank. "FundsChain" feature / World Bank blockchain initiatives (2025).
\item World Economic Forum. "Inclusive Deployment of Blockchain for Supply Chains" (2019).
\item GS1 US. "Applying GS1 Standards for Supply Chain Visibility in Blockchain" (2020, EPCIS guidance).
\item NIST. "Manufacturing Supply Chain Traceability with Blockchain Related Technology" (2023 reference implementation).
\item Hyperledger Fabric documentation (IBM / Hyperledger community docs).
\item Malaysia Securities Commission. "Guidelines on Islamic Capital Market Products and Services" (contains musharakah/mudharabah guidance).
\item AAOIFI standards and IFSB guidance on Islamic finance products.
\item Deloitte. "Using blockchain to drive supply chain transparency" (industry article / whitepaper).
\item Oracle \& Deloitte. "Value-driven supply chain: blockchain + IoT" (POV briefing).
\item IBM Institute for Business Value — Blockchain for sustainable supply chains whitepaper.
\item World Food Programme — "Building Blocks" (identity and cash transfer on blockchain) case study.
\item Citrusxchange / SupplyChainDigital coverage on invoice-financing platforms.
\item Academic literature: Govindan et al. (2023), Rakshit et al. (2022), Alimohammadlou et al. (2023), Ojo \& Zhou (2020), Pi (2026) — supply chain and blockchain adoption papers.
\item WEF publications on government procurement \& blockchain (IDB collaboration case studies).
\item ISO 20022 payment messaging guidance (SWIFT whitepapers).
\item W3C DID / Verifiable Credentials specifications.
\item White \& Case LLP: "Tokenised Islamic finance products" insights.
\item R3 Corda / Enterprise Ethereum (Besu/Quorum) platform docs (privacy \& permissioning comparisons).
\item Various GitHub POCs (IBM Food Trust sample chaincode, MangoChain, public repos demonstrating invoice tokenisation).
\item Industry articles: ISM Supply Management, SupplyChainDigital.
\item McKinsey, PwC industry reports on blockchain in supply chains (selected briefs).
\item National fintech regulator guidance (e.g., Bahrain, UAE tokenisation frameworks where public).
\item FAO / UNECE reports on blockchain for agriculture \& market access.
\item UN / IDB procurement pilot summaries (public reports).
\item Research articles on Islamic fintech and blockchain (Fintech in Islamic finance literature review — 2022/2023).
\end{enumerate}

\section{Appendix B: Next steps and recommended actions}
\begin{enumerate}
\item Validate high-risk requirements with named stakeholders (Shariah board, pilot bank, procurement owner).\
\item Prioritize requirements by legal/compliance risk and business value; produce top-10 MVP backlog.\
\item Produce sequence diagrams and user-stories for top 8 requirements.\
\item Build prototype: identity + PO lifecycle + simple PLS contract + tokenised invoice flow in a sandbox.
\end{enumerate}

\section{Appendix C: Assumptions}
\begin{itemize}
\item Permissioned blockchain baseline (Fabric/Besu/Corda-like).\
\item Primary jurisdiction not yet selected; retention rules configurable.\
\item Stakeholder roles are placeholders for later mapping to real organizations.\
\end{itemize}

\vfill\noindent
\hrulefill\\
\noindent {\footnotesize This preliminary SRS is a public-source-derived draft and requires stakeholder validation.}
\end{document}
